\section{Exact factor exchange inside a residual fiber}
\label{sec:tpc57-factor-exchange}

We work only with the literal fixed-shift carrier inherited from
TPC--55--TPC--56.  Fix \(h_0\ne0\), put
\[
 H:=\operatorname{rad}|h_0|,
\]
and take \(X\) sufficiently large that \(L>|h_0|\).  A terminal atom is
written
\[
 u=(d\ell,r,j,\gamma),
 \qquad
 d\in[D,2D],\quad \ell\in[L,2L]\cap\mathbb P,
\]
where \(d,r,dr\) are squarefree, \((d,r)=1\), and
\begin{equation}
 rp-d\ell j=h_0
 \label{eq:tpc57-terminal-equation}
\end{equation}
for the terminal prime \(p\), with every retained-row, cofactor, residue,
profile, and terminal-support condition imposed exactly as in the literal
carrier.  Its coarse output key is
\[
 q(u)=(\gamma,s,j),\qquad s=dr.
\]
No ambient Diophantine solution is called an atom unless it belongs to
that exact carrier.

\subsection{The factor-swap normal form}

\begin{proposition}[Unique factor exchange]
\label{prop:tpc57-factor-exchange}
Let
\[
 u_i=(d_i\ell_i,r_i,j,\gamma),\qquad i=1,2,
\]
be two terminal atoms in the same coarse output fiber.  There are unique
positive, pairwise coprime, squarefree integers \(a,b,c,e\) such that
\begin{equation}
 \boxed{
 d_1=ab,\qquad d_2=ac,\qquad
 r_1=ce,\qquad r_2=be,
 \qquad s=abce.}
 \label{eq:tpc57-factor-swap}
\end{equation}
They are given by
\begin{equation}
 a=(d_1,d_2),\qquad
 b=d_1/a,\qquad c=d_2/a,
 \qquad e=s/(abc)=(r_1,r_2).
 \label{eq:tpc57-factor-swap-parameters}
\end{equation}
Conversely, any tuple of the form
\eqref{eq:tpc57-factor-swap} that satisfies both literal terminal
equations and all carrier restrictions gives two atoms in one coarse
fiber.
\end{proposition}

\begin{proof}
Because \(s\) is squarefree and \(d_i r_i=s\) with
\((d_i,r_i)=1\), every prime of \(s\) lies in exactly one of four
disjoint classes: in both \(d_1,d_2\), in \(d_1,r_2\), in
\(d_2,r_1\), or in both \(r_1,r_2\).  Their products are respectively
\(a,b,c,e\), which proves existence, pairwise coprimality, and
\eqref{eq:tpc57-factor-swap}.  The four prime classes are recovered from
\((d_1,d_2,r_1,r_2)\), so the representation is unique.  The converse is
immediate after the literal support conditions, rather than only the
displayed equations, are imposed.
\end{proof}

The normal form may also be viewed relative to one atom.  If
\(d_1=d\) and \(r_1=r\), a different divisor choice is obtained by
choosing \(b\mid d\) and \(c\mid r\) and exchanging these two coprime
factor packets:
\begin{equation}
 d_2=\frac d b c,
 \qquad
 r_2=\frac r c b.
 \label{eq:tpc57-relative-factor-swap}
\end{equation}
The pair \((b,c)\) is unique once the partner divisor \(d_2\) is known.

\begin{proposition}[Exchange identity and exact divisibility]
\label{prop:tpc57-exchange-identity}
Under the hypotheses of
\cref{prop:tpc57-factor-exchange}, the two terminal equations imply
\begin{equation}
 \boxed{
 e(cp_1-bp_2)=aj(b\ell_1-c\ell_2).}
 \label{eq:tpc57-exchange-identity}
\end{equation}
Moreover, for \(i=1,2\),
\begin{equation}
 (r_i,j)=(r_i,h_0),
 \label{eq:tpc57-terminal-gcd-in-factor-section}
\end{equation}
and consequently
\begin{equation}
 \boxed{
 \frac{e}{(e,h_0)}\mid b\ell_1-c\ell_2.}
 \label{eq:tpc57-exchange-divisibility}
\end{equation}
\end{proposition}

\begin{proof}
Substitute \eqref{eq:tpc57-factor-swap} into the two copies of
\eqref{eq:tpc57-terminal-equation} and subtract:
\[
 cep_1-ab\ell_1j= bep_2-ac\ell_2j.
\]
This is \eqref{eq:tpc57-exchange-identity}.

For sufficiently large \(X\), one has \((\ell_i,r_i)=1\).  Indeed,
the prime \(\ell_i\) dividing \(r_i\) would divide both terms on the
left of \eqref{eq:tpc57-terminal-equation}, and hence would divide the
fixed nonzero integer \(h_0\), contrary to \(\ell_i\ge L>|h_0|\).
Thus \((d_i\ell_i,r_i)=1\), and reduction of the terminal equation
modulo \(r_i\) gives
\[
 j\equiv-h_0(d_i\ell_i)^{-1}\pmod {r_i}.
\]
Multiplication by a unit modulo \(r_i\) proves
\eqref{eq:tpc57-terminal-gcd-in-factor-section}.

Finally \((a,e)=1\), so \eqref{eq:tpc57-exchange-identity} gives
\(e\mid j(b\ell_1-c\ell_2)\).  Since
\((e,j)=(e,h_0)\), division by this gcd proves
\eqref{eq:tpc57-exchange-divisibility}.
\end{proof}

\subsection{Same-divisor and cross-divisor separation}

\begin{corollary}[Same-divisor collisions require a short cofactor]
\label{cor:tpc57-same-divisor-short}
Suppose \(u_1\ne u_2\) and \(d_1=d_2=d\).  Then \(r_1=r_2=r\),
\(\ell_1\ne\ell_2\), and
\begin{equation}
 \frac r{(r,h_0)}\mid \ell_1-\ell_2.
 \label{eq:tpc57-same-divisor-spacing}
\end{equation}
In particular,
\begin{equation}
 \boxed{r\le HL.}
 \label{eq:tpc57-same-divisor-bound}
\end{equation}
Therefore, on a cofactor block \(r\in[R,2R]\) with \(R>HL\), a fixed
divisor \(d\) contributes at most one terminal atom to any coarse output
fiber.
\end{corollary}

\begin{proof}
In the unique factorization of
\cref{prop:tpc57-factor-exchange}, equality \(d_1=d_2\) gives
\(b=c=1\), \(a=d\), and \(e=r\).  Distinctness then forces
\(\ell_1\ne\ell_2\), while
\eqref{eq:tpc57-exchange-divisibility} gives
\eqref{eq:tpc57-same-divisor-spacing}.  Since \(r\) is squarefree,
\((r,h_0)\mid H\), and since two points of \([L,2L]\) differ by at most
\(L\), the asserted bound follows.
\end{proof}

\begin{proposition}[Large exchanged factors in a balanced block]
\label{prop:tpc57-large-exchange}
Assume in addition that \(L>4D\), that \(d_1\ne d_2\), and that
\begin{equation}
 d_i\in[D,2D],\qquad r_i\in[R,2R]\qquad(i=1,2).
 \label{eq:tpc57-common-dyadic-block}
\end{equation}
Then
\begin{equation}
 b\ell_1-c\ell_2\ne0
 \label{eq:tpc57-nonzero-source-exchange}
\end{equation}
and the exchanged factors in
\eqref{eq:tpc57-factor-swap} obey
\begin{equation}
 \boxed{
 b^2\ge\frac{R}{6HL},
 \qquad
 c^2\ge\frac{R}{6HL}.}
 \label{eq:tpc57-large-exchange-bound}
\end{equation}
Equivalently, both exchanged factor packets have size at least
\[
 \left(\frac{R}{6HL}\right)^{1/2}.
\]
\end{proposition}

\begin{proof}
If \(b\ell_1=c\ell_2\), pairwise coprimality of \(b,c\) implies
\(b\mid\ell_2\) and \(c\mid\ell_1\).  But
\(b,c\le2D<L/2\le\ell_i/2\), and the \(\ell_i\) are prime.  Hence
\(b=c=1\), contrary to \(d_1\ne d_2\).  This proves
\eqref{eq:tpc57-nonzero-source-exchange}.

By \eqref{eq:tpc57-exchange-divisibility}, squarefreeness of \(e\), and
the definition of \(H\),
\[
 e\le H|b\ell_1-c\ell_2|
 \le2HL(b+c).
\]
The ratio \(b/c=d_1/d_2\) lies in \([1/2,2]\).  Thus
\(e\le6HLb\) and, symmetrically, \(e\le6HLc\).  Since
\(r_2=be\ge R\) and \(r_1=ce\ge R\), one also has
\(e\ge R/b\) and \(e\ge R/c\).  Combining the two pairs of inequalities
gives \eqref{eq:tpc57-large-exchange-bound}.
\end{proof}

\subsection{The exact partner candidate set}

Fix a terminal atom \(u=(d\ell,r,j,\gamma)\).  Define
\(\mathfrak C(u)\) to be the set of triples \((b,c,\ell')\) for which
\begin{enumerate}[label=\textup{(\roman*)},leftmargin=3em]
 \item \(b\mid d\), \(c\mid r\), and
       \((b,c)\ne(1,1)\);
 \item
       \[
        d'=\frac d b c,
        \qquad r'=\frac r c b
       \]
       are integers in the declared divisor and cofactor supports;
 \item \(\ell'\) is in the exact declared source-prime support and the
       atom
       \[
        u'=\left(d'\ell',r',j,\gamma\right)
       \]
       belongs to the literal terminal-admissible universe, including all
       retained-pair, residue, cutoff, and carrier restrictions.
\end{enumerate}
The terminal prime of \(u'\) is not a free coordinate: it is exactly
\begin{equation}
 p'=\frac{d'\ell'j+h_0}{r'}.
 \label{eq:tpc57-partner-prime}
\end{equation}

\begin{proposition}[Exact partner parameterization]
\label{prop:tpc57-exact-partner-set}
The map
\begin{equation}
 u'\longmapsto
 \left(\frac d{(d,d')},
       \frac{d'}{(d,d')},\ell'\right)
 \label{eq:tpc57-partner-map}
\end{equation}
is a bijection from the distinct-divisor terminal partners of \(u\) in
its full coarse output fiber onto \(\mathfrak C(u)\).  Same-divisor
partners are exactly the further source primes on the single determinant
ladder with \(d'=d\) and \(r'=r\).
\end{proposition}

\begin{proof}
For a partner \(u'\), apply
\cref{prop:tpc57-factor-exchange} with \(d_1=d,r_1=r\).  It gives the
unique divisors \(b=d/(d,d')\) and \(c=d'/(d,d')\), and
\eqref{eq:tpc57-relative-factor-swap}; literal membership gives all the
conditions defining \(\mathfrak C(u)\).  Conversely, every element of
\(\mathfrak C(u)\) defines a literal terminal atom with
\(d'r'=dr=s\), the same \(j,\gamma\), and hence the same output key.
Uniqueness of the factor swap and of the source row proves bijectivity.
The case \(d'=d\) has \(b=c=1\) and is precisely the fixed-divisor
ladder.
\end{proof}

Assume now that the root cofactor satisfies \(r\in[R,2R]\), and constrain
its partners to the same block \(r'\in[R,2R]\).  Define
\begin{equation}
 \mathfrak E_R(d,r)
 :=\left\{(b,c):\exists\ell'\;
   \begin{array}{l}
   (b,c,\ell')\in\mathfrak C(u),\\[-2pt]
   r'=(r/c)b\in[R,2R]
   \end{array}\right\}.
 \label{eq:tpc57-same-block-exchange-set}
\end{equation}
The fixed-divisor ladder spacing gives the
deterministic degree estimate
\begin{equation}
 \#\{u'\ne u:q(u')=q(u),\ r'\in[R,2R]\}
 \le
 \left\lfloor\frac{HL}{R}\right\rfloor
 +\left(1+\left\lfloor\frac{HL}{R}\right\rfloor\right)
   |\mathfrak E_R(d,r)|.
 \label{eq:tpc57-partner-degree-bound}
\end{equation}
When \(d'\ne d\), every pair counted by
\(\mathfrak E_R(d,r)\) additionally satisfies
\eqref{eq:tpc57-large-exchange-bound}.  This is sharper structural
information than the untruncated divisor bound
\(|\mathfrak E_R(d,r)|\le\tau(d)\tau(r)\), although it does not by itself
force the candidate set to be empty.

\begin{corollary}[High-cofactor structural exclusion]
\label{cor:tpc57-high-cofactor-exclusion}
Under the hypotheses above, if
\begin{equation}
 R>24HLD^2,
 \label{eq:tpc57-high-cofactor-threshold}
\end{equation}
then every coarse terminal fiber restricted to
\(d\in[D,2D]\) and \(r\in[R,2R]\) contains at most one atom.
\end{corollary}

\begin{proof}
The hypothesis implies \(R>HL\), so
\cref{cor:tpc57-same-divisor-short} excludes a same-divisor pair.  For a
cross-divisor pair, \(b\le d_1\le2D\), whereas
\eqref{eq:tpc57-large-exchange-bound} and
\eqref{eq:tpc57-high-cofactor-threshold} give \(b^2>4D^2\), a
contradiction.
\end{proof}

There is also a full-fiber version, with only a harmless change of the
constant.  It is useful when a core cofactor block is used merely to find
energy and no block-restricted output is introduced.

\begin{corollary}[A high-cofactor atom is a full-fiber singleton]
\label{cor:tpc57-high-cofactor-full-singleton}
Assume \(L>4D\).  Let \(u_1\) be a terminal atom with \(d_1\in[D,2D]\) and
\(r_1\in[R,2R]\), while every possible source divisor in its full literal
carrier remains in \([D,2D]\).  If
\begin{equation}
 R>48HLD^2,
 \label{eq:tpc57-full-singleton-threshold}
\end{equation}
then \(u_1\) has no distinct terminal partner anywhere in its full coarse
output fiber.
\end{corollary}

\begin{proof}
Every partner has
\(r_2/r_1=d_1/d_2\in[1/2,2]\), hence
\(r_2\in[R/2,4R]\).  A same-divisor partner is excluded by
\cref{cor:tpc57-same-divisor-short}.  For a cross-divisor partner, the
proof of \cref{prop:tpc57-large-exchange} still gives
\(e\le6HLb\), while \(r_2=be\ge R/2\).  Therefore
\[
 b^2\ge\frac{R}{12HL}>4D^2,
\]
contradicting \(b\le d_1\le2D\).
\end{proof}

The last corollary is a support theorem only.  In fact, the inherited
large-prime carrier has a global cofactor ceiling that is disjoint from
this high-cofactor regime; see
\cref{eq:tpc57-large-prime-r-ceiling,eq:tpc57-high-cofactor-global-gap}.
Thus the corollary is a sharp structural boundary, not an activation
mechanism for the present physical packet.

Finally, the factor geometry localizes possible partners up to a constant
scale enlargement.  If one atom has \(r_1\in[R,2R]\), then every atom in
its full coarse fiber satisfies
\begin{equation}
 \frac12\le\frac{r_2}{r_1}=\frac{d_1}{d_2}\le2,
 \qquad
 r_2\in[R/2,4R].
 \label{eq:tpc57-cofactor-saturation}
\end{equation}
Thus a singleton audit for a core block needs to inspect only its
constant-factor saturation, not remote cofactor scales.  This observation
does not license a lower bound for a block-restricted output to be returned
to the full sum; it merely identifies the complete set in which a full
fiber partner can occur.
